% Part: first-order-logic
% Chapter: natural-deduction
% Section: derivations

\documentclass[../../../include/open-logic-section]{subfiles}

\begin{document}

\iftag{FOL}
      {\olfileid{fol}{ntd}{der}}
      {\olfileid{pl}{ntd}{der}}

\olsection{\usetoken{P}{derivation}}

\begin{explain}
We weten nu wat een aanname is en welke afleidingsregels er zijn.
Daarmee bouwen we !!^{derivation}s in natuurlijke deductie stap voor
stap op. Elke !!{derivation} is ofwel één losse aanname, of bestaat uit
één, twee of drie eerdere !!{derivation}s met daaronder een
correcte afleidingsstap.
\end{explain}

\begin{defn}[!!^{derivation}]
\Article{derivation} \emph{!!{derivation}} van !!a{sentence}~$!A$ uit
aannamen~$\Gamma$ is een eindige boom van !!{sentence}s. Die boom moet
aan de volgende voorwaarden voldoen:
\begin{enumerate}
\item Elke bovenste !!{sentence} behoort tot $\Gamma$ of wordt door een
  afleidingsstap in de boom !!{discharged}.
\item De onderste !!{sentence} is~$!A$.
\item Elke andere !!{sentence} is een premisse van een correct
  toegepaste afleidingsregel. De conclusie van die regel staat direct
  onder die !!{sentence} in de boom. De enige uitzondering is de
  onderste zin~$!A$.
\end{enumerate}
De onderste zin $!A$ heet dan de \emph{conclusie} van de
!!{derivation}. De verzameling $\Gamma$ bestaat uit de
!!{undischarged} aannamen ervan.

Bestaat er !!a{derivation} van $!A$ uit~$\Gamma$, dan is $!A$ uit
$\Gamma$ \emph{!!{derivable}}. In symbolen schrijven we dat als:
$\Gamma \Proves !A$. Bestaat er !!a{derivation} van~$!A$ waarin elke
aanname is !!{discharged}, dan schrijven we~$\Proves !A$.
\end{defn}

\begin{ex}
Elke losse aanname is op zichzelf !!a{derivation}. Daarom is $!A$
!!a{derivation}, en $!B$ ook. We zetten die twee afleidingen naast
elkaar en passen er de regel $\Intro{\land}$ op toe. Zo krijgen we een
nieuwe !!{derivation}:
\begin{prooftree}
\AxiomC{$!A$}
\AxiomC{$!B$}
\RightLabel{\Intro{\land}}
\BinaryInfC{$!A \land !B$}
\end{prooftree}
De regels zijn algemeen. We mogen de $!A$ en~$!B$ door willekeurige
!!{sentence}s vervangen, bijvoorbeeld door $!C$ en~$!D$. De conclusie
wordt dan $!C \land !D$. Daarom is
\begin{prooftree}
\AxiomC{$!C$}
\AxiomC{$!D$}
\RightLabel{\Intro{\land}}
\BinaryInfC{$!C \land !D$}
\end{prooftree}
een correcte !!{derivation}. We mogen de rollen ook omdraaien: laat
$!D$ de rol van~$!A$ spelen en $!C$ die van~$!B$. Dan is
\begin{prooftree}
\AxiomC{$!D$}
\AxiomC{$!C$}
\RightLabel{\Intro{\land}}
\BinaryInfC{$!D \land !C$}
\end{prooftree}
eveneens een correcte !!{derivation}.

Vervolgens kunnen we bijvoorbeeld de regel $\Intro{\lif}$ toepassen.
Die regel vormt een implicatie en laat ons elke aanname !!{discharge}
die gelijk is aan het antecedent van die implicatie. Dat levert de
volgende twee correcte !!{derivation}s op:
\begin{prooftree}
\AxiomC{$\Discharge{!C}{1}$}
\AxiomC{$!D$}
\RightLabel{\Intro{\land}}
\BinaryInfC{$!C \land !D$}
\DischargeRule{\Intro{\lif}}{1}
\UnaryInfC{$!C \lif (!C \land !D)$}
\DisplayProof\bottomAlignProof
\AxiomC{$!C$}
\AxiomC{$\Discharge{!D}{1}$}
\RightLabel{\Intro{\land}}
\BinaryInfC{$!C \land !D$}
\DischargeRule{\Intro{\lif}}{1}
\UnaryInfC{$!D \lif (!C \land !D)$}
\end{prooftree}
De eerste toont aan dat $!D \Proves !C \lif (!C \land !D)$; de tweede
toont aan dat $!C \Proves !D \lif (!C \land !D)$.

Het intrekken van aannamen is toegestaan, maar niet verplicht. We
hoeven een aanname dus niet in te trekken. Een regel mag zelfs worden
toegepast als de bedoelde aanname niet in de !!{derivation} voorkomt.
Het volgende is bijvoorbeeld toegestaan, hoewel er geen aanname~$!A$
is om te worden !!{discharged}:
\begin{prooftree}
  \AxiomC{$!B$}
  \DischargeRule{\Intro{\lif}}{1}
  \UnaryInfC{$!A \lif !B$}
\end{prooftree}
\end{ex}
\end{document}
